\chapter{Efficiency}

\textbf{Readings:} none in CP:AMA (this material is not in the textbook).

\section{Algorithms and Efficiency}

\begin{definition}
An \emph{algorithm} is a step-by-step description of how to solve a problem.
\end{definition}

Different algorithms for the same problem can take very different amounts of resources.
\begin{itemize}
  \item \textbf{Time efficiency:} how long it takes an algorithm to solve a problem (the most
    common measure).
  \item \textbf{Space efficiency:} how much memory it requires.
\end{itemize}
Unless stated otherwise, ``efficiency'' means time efficiency.

\section{Running Time}

Measuring with a stopwatch depends on the machine, the compiler and the input, so instead we
count the number of \emph{elementary operations} (arithmetic, comparisons, assignments, array
indexing, function calls, \dots), each of which takes constant time.

The running time depends on the size of the input, so it is a function $T(n)$, where $n$ is the
size of the input. What $n$ means depends on the function: the length of an array, the number of
nodes in a tree, the value of an integer argument, the length of a string. Other letters (such
as $m$ or $k$) are common when there are several inputs; always state what each variable means.

\subsection{Worst case}

The running time can depend on the \emph{contents} of the input, not just its size (e.g.\ linear
search finds the item immediately if it is first). We usually care about the \textbf{worst-case
running time}: the maximum over all inputs of size $n$. It is a conservative guarantee: the
function will never be slower than this. (Best case is rarely useful; average case is useful but
harder to analyse.)

\section{Big-O Notation}

In practice we are interested in the \emph{order} (growth rate) of the running time, not the
exact constants. For example $17n + 1$ and $8n + 7$ are both ``order $n$'', written $\bigO(n)$.

\begin{definition}
$f(n) \in \bigO(g(n))$ if there exist constants $c > 0$ and $n_0 > 0$ such that
$f(n) \le c \cdot g(n)$ for all $n \ge n_0$.
\end{definition}

\begin{example}
$17n + 1 \in \bigO(n)$: take $c = 18$, $n_0 = 1$; then $17n + 1 \le 18n$ for all $n \ge 1$.
Also $17n + 1 \in \bigO(n^2)$, but that is not a useful (tight) bound. In CS 136 we always give
the tightest simple bound.
\end{example}

\subsection{Rules for simplifying}
\begin{itemize}
  \item Drop constant factors: $\bigO(5n) = \bigO(n)$.
  \item Keep only the dominant term: $\bigO(3n^2 + 100n + 7) = \bigO(n^2)$.
  \item Sum: $\bigO(f) + \bigO(g) = \bigO(\max(f, g))$ (sequential code).
  \item Product: $\bigO(f) \cdot \bigO(g) = \bigO(f \cdot g)$ (nested loops).
  \item Bases of logarithms do not matter: $\log_2 n$ and $\log_{10} n$ differ by a constant
    factor, so we just write $\bigO(\log n)$.
\end{itemize}

\subsection{Common orders}

\begin{center}
\begin{tabular}{lll}
  \toprule
  Order & Name & Example \\
  \midrule
  $\bigO(1)$ & constant & array indexing, \code{first} \\
  $\bigO(\log n)$ & logarithmic & binary search \\
  $\bigO(n)$ & linear & linear search, \code{strlen} \\
  $\bigO(n \log n)$ & linearithmic & merge sort \\
  $\bigO(n^2)$ & quadratic & selection sort \\
  $\bigO(n^3)$ & cubic & naive matrix multiplication \\
  $\bigO(2^n)$ & exponential & naive Fibonacci \\
  \bottomrule
\end{tabular}
\end{center}
Each is strictly worse than the one above it for large $n$. For $n = 10^6$: $\log n \approx 20$,
$n \log n \approx 2\times 10^7$ (instant), $n^2 = 10^{12}$ (many minutes), and $2^n$ is
hopeless.

\begin{aside}
$\bigO$ is an \emph{upper} bound. CS 240 introduces $\Omega$ (lower bound) and $\Theta$ (tight
bound). When CS 136 says ``the running time is $\bigO(n)$'' it really means $\Theta(n)$.
\end{aside}

\section{Analysing Iterative Code}

\begin{enumerate}
  \item Each elementary statement is $\bigO(1)$.
  \item A sequence of blocks takes the sum (i.e.\ the maximum) of their times.
  \item A loop takes (number of iterations) $\times$ (time per iteration), or, if the time per
    iteration varies, the sum over all iterations.
  \item A function call takes the running time of the function (not $\bigO(1)$!).
\end{enumerate}

\begin{example}
\begin{lstlisting}[style=c]
for (int i = 0; i < n; ++i) {          // n iterations
  for (int j = 0; j < n; ++j) {        //   n iterations each
    sum += i * j;                      //     O(1)
  }
}
\end{lstlisting}
$n \cdot n \cdot \bigO(1) = \bigO(n^2)$.
\end{example}

\begin{example}[Selection sort]
The inner loop runs $n - 1 - i$ times for $i = 0, \dots, n-2$, so the total is
\[ \sum_{i=0}^{n-2} (n - 1 - i) = (n-1) + (n-2) + \dots + 1 = \frac{n(n-1)}{2} \in \bigO(n^2). \]
Selection sort is $\bigO(n^2)$ in \emph{every} case.
\end{example}

\begin{example}[Halving loop]
\begin{lstlisting}[style=c]
while (n > 1) {
  n /= 2;
  ++count;
}
\end{lstlisting}
$n$ is halved each iteration, so there are about $\log_2 n$ iterations: $\bigO(\log n)$. This is
why binary search is $\bigO(\log n)$.
\end{example}

\begin{important}[Hidden costs]
Watch out for function calls in loop conditions:
\begin{lstlisting}[style=c]
for (int i = 0; i < strlen(s); ++i) { ... }   // O(n^2)! strlen is O(n)
\end{lstlisting}
Compute \code{strlen(s)} once before the loop to make it $\bigO(n)$.
\end{important}

\section{Analysing Recursive Code}

For a recursive function, write a \emph{recurrence relation} for $T(n)$: the time for the
non-recursive work plus the time of the recursive calls. Then look it up in the table (you are
expected to know these; deriving them is covered in later courses).

\begin{center}
\begin{tabular}{ll}
  \toprule
  Recurrence & Solution \\
  \midrule
  $T(n) = \bigO(1) + T(n-k)$ & $\bigO(n)$ \\
  $T(n) = \bigO(n) + T(n-k)$ & $\bigO(n^2)$ \\
  $T(n) = \bigO(n^2) + T(n-k)$ & $\bigO(n^3)$ \\
  $T(n) = \bigO(1) + T(n/k)$ & $\bigO(\log n)$ \\
  $T(n) = \bigO(1) + k \cdot T(n/k)$ & $\bigO(n)$ \\
  $T(n) = \bigO(n) + k \cdot T(n/k)$ & $\bigO(n \log n)$ \\
  $T(n) = \bigO(1) + T(n-k_1) + T(n-k_2)$ & $\bigO(2^n)$ \\
  \bottomrule
\end{tabular}
\end{center}
(for constants $k, k_1, k_2 \ge 1$; for the $n/k$ rows, $k \ge 2$).

\begin{example}
\begin{itemize}
  \item \code{sum\_first(n)} (Unit 02) does $\bigO(1)$ work and recurses on $n-1$: $\bigO(n)$.
  \item Recursive binary search: $T(n) = \bigO(1) + T(n/2)$: $\bigO(\log n)$.
  \item Naive \code{fib(n) = fib(n-1) + fib(n-2)}: $\bigO(2^n)$.
  \item A Racket function that calls \code{append} or \code{length} ($\bigO(n)$) on the rest of
    the list at each step: $T(n) = \bigO(n) + T(n-1) = \bigO(n^2)$.
\end{itemize}
\end{example}

\section{Sorting Algorithms Revisited}

\subsection{Insertion sort}

In the worst case (reverse-sorted input) inserting element $i$ takes $i$ swaps, so the total is
$1 + 2 + \dots + (n-1) \in \bigO(n^2)$. In the best case (already sorted) each insertion is
$\bigO(1)$, giving $\bigO(n)$.

\subsection{Quicksort}

Partitioning is $\bigO(n)$. If the pivot always splits the array in half,
$T(n) = \bigO(n) + 2T(n/2) = \bigO(n \log n)$. But in the worst case (e.g.\ an already sorted array
with the first element as pivot) one side is empty every time: $T(n) = \bigO(n) + T(n-1) =
\bigO(n^2)$. Quicksort's \emph{average} case is $\bigO(n \log n)$, and it is very fast in practice,
especially with a random pivot.

\subsection{Merge sort}

Merge sort guarantees $\bigO(n \log n)$: split the array into halves, sort each half recursively,
and \emph{merge} the two sorted halves in $\bigO(n)$ time.
\begin{lstlisting}[style=c]
// merge(dest, src1, len1, src2, len2) merges the sorted arrays src1 and
//   src2 into dest
// requires: dest has length len1 + len2 and does not overlap src1/src2
void merge(int dest[], const int src1[], int len1,
           const int src2[], int len2) {
  int i = 0, j = 0, k = 0;
  while (i < len1 && j < len2) {
    if (src1[i] <= src2[j]) {
      dest[k++] = src1[i++];
    } else {
      dest[k++] = src2[j++];
    }
  }
  while (i < len1) dest[k++] = src1[i++];
  while (j < len2) dest[k++] = src2[j++];
}

void merge_sort(int a[], int len) {
  if (len <= 1) return;
  int llen = len / 2;
  int rlen = len - llen;
  int *left = malloc(llen * sizeof(int));    // see Unit 10
  int *right = malloc(rlen * sizeof(int));
  for (int i = 0; i < llen; ++i) left[i] = a[i];
  for (int i = 0; i < rlen; ++i) right[i] = a[llen + i];
  merge_sort(left, llen);
  merge_sort(right, rlen);
  merge(a, left, llen, right, rlen);
  free(left);
  free(right);
}
\end{lstlisting}
$T(n) = \bigO(n) + 2T(n/2) = \bigO(n \log n)$ in every case, at the cost of $\bigO(n)$ extra space.

\begin{center}
\begin{tabular}{lccc}
  \toprule
  Algorithm & Best & Worst & Extra space \\
  \midrule
  Selection sort & $\bigO(n^2)$ & $\bigO(n^2)$ & $\bigO(1)$ \\
  Insertion sort & $\bigO(n)$ & $\bigO(n^2)$ & $\bigO(1)$ \\
  Quicksort & $\bigO(n \log n)$ & $\bigO(n^2)$ & $\bigO(\log n)$ to $\bigO(n)$ stack \\
  Merge sort & $\bigO(n \log n)$ & $\bigO(n \log n)$ & $\bigO(n)$ \\
  \bottomrule
\end{tabular}
\end{center}

\section{Running Times of Library Functions}

\subsection{Racket}
\begin{itemize}
  \item $\bigO(1)$: \code{cons}, \code{cons?}, \code{empty}, \code{empty?}, \code{first},
    \code{rest}, \code{second}, arithmetic on small numbers, \code{equal?} on simple values.
  \item $\bigO(n)$: \code{length}, \code{last}, \code{reverse}, \code{append} (in the length of
    the first list), \code{list-ref} (in the index), \code{member}.
  \item Abstract list functions (\code{map}, \code{filter}, \code{foldr}) are $\bigO(n)$ when the
    function they apply is $\bigO(1)$.
  \item Racket's \code{sort} is $\bigO(n \log n)$.
\end{itemize}

\subsection{C}
\begin{itemize}
  \item All the \code{<string.h>} functions are $\bigO(n)$, where $n$ is the length of the string.
    For \code{strcmp}, $n$ is the length of the \emph{shorter} string; for \code{strcat},
    \code{strlen(dest) + strlen(src)}.
  \item \code{printf} and \code{scanf} are $\bigO(1)$ for numbers and characters, and $\bigO(n)$
    for strings (\code{\%s}) of length $n$.
  \item Array indexing and pointer arithmetic are $\bigO(1)$.
  \item \code{malloc} and \code{free} are $\bigO(1)$ for our purposes; \code{realloc} is $\bigO(n)$
    where $n$ is the size of the block (Unit 10).
\end{itemize}

\section{Documenting Efficiency}

From now on, the documentation of every function includes its time efficiency, with the meaning
of each variable:
\begin{lstlisting}[style=c]
// find_sorted(item, a, len) produces the index of item in a, or -1
// requires: a is sorted in increasing order
// time: O(log n), where n is len
\end{lstlisting}

\section{Space Efficiency}

Space efficiency is analysed the same way: how much additional memory does the function need as
a function of $n$? Remember that recursion uses stack space: a non-tail recursive function with
depth $n$ uses $\bigO(n)$ space even if it allocates nothing, whereas the equivalent loop uses
$\bigO(1)$.

\section{Summary}
\begin{itemize}
  \item Measure efficiency by counting elementary operations as a function of input size $n$;
    focus on the worst case.
  \item Big-O ignores constant factors and lower-order terms; know the common orders and their
    ranking.
  \item Loops: iterations $\times$ cost; nested loops multiply; watch hidden $\bigO(n)$ calls.
  \item Recursion: write a recurrence and use the table.
  \item Selection sort $\bigO(n^2)$; insertion sort $\bigO(n^2)$ worst, $\bigO(n)$ best; quicksort
    $\bigO(n \log n)$ average, $\bigO(n^2)$ worst; merge sort $\bigO(n \log n)$ always.
  \item Document \emph{time:} in every contract.
\end{itemize}
